% answers_short.tex —— 2009 年·普通高等学校招生考试·北京卷·文科 · 简短答案速查
% 由 tools/make_short_answers.py 自 answers.tex 的 \gksolhead/\gkanswer 提取，
% 并与 2009/2009_数学真题_北京文_答案要点.md 逐题比对后生成，勿手改（改源 MD 后重跑）。

\begin{gknotice}
\textbf{使用说明}：本节只给最终答案便于快速核对，完整推导见解析版。答案为编写者独立推导，\textbf{非官方标准答案}。
\end{gknotice}

\gkpart{选择题}
\begin{center}\begin{tabular}{ccccccccc}
\toprule
题号 & 1 & 2 & 3 & 4 & 5 & 6 & 7 & 8 \\
\midrule
答案 & A & D & B & C & C & A & D & D \\
\bottomrule
\end{tabular}\end{center}

\gkpart{填空题}
\begin{itemize}[leftmargin=2.6em]
  \item[\textbf{9.}] $-\frac35$。
  \item[\textbf{10.}] $16$；$255$。
  \item[\textbf{11.}] $9$。
  \item[\textbf{12.}] $\log_3 2$。
  \item[\textbf{13.}] $2$；$\frac{2\pi}{3}$。
  \item[\textbf{14.}] 0（按原页字面）；官方同题为 6。
\end{itemize}

\gkpart{解答题}
\begin{itemize}[leftmargin=2.6em]
  \item[\textbf{15.}] （1）$\pi$；（2）最大值 $1$，最小值 $-\frac{\sqrt3}{2}$。
  \item[\textbf{16.}] （1）证明见解析；（2）$45^\circ$。
  \item[\textbf{17.}] （1）$\frac4{27}$；（2）$\frac89$。
  \item[\textbf{18.}] （1）$a=4$，$b=24$；（2）增区间 $(-\infty,-2)$、$(2,+\infty)$，减区间 $(-2,2)$；极大值 $f(-2)=40$，极小值 $f(2)=8$。
  \item[\textbf{19.}] （1）$x^2-\frac{y^2}{2}=1$；（2）$m=\pm1$。
  \item[\textbf{20.}] （1）$b_3=7$；（2）$m^2+2m$；（3）存在，$p=\frac13$，$q\in\left[-\frac23,-\frac13\right)$。
\end{itemize}
